% ---
% Inicia os anexos
% ---
\begin{anexosenv}
    \chapter{Código Exercício 1.6}\label{cap:ex1-6-codigo}
    
    \begin{lstlisting}[language=Python, breaklines=true, basicstyle=\tiny]
    import numpy as np
    import matplotlib.pyplot as plt
    import scipy.special

    # --- constantes ----
    CENTRAL_FLUX = 0
    UPWIND_FLUX = 1

    A = np.array([0.0,
    -567301805773.0 / 1357537059087.0,
    -2404267990393.0 / 2016746695238.0,
    -3550918686646.0 / 2091501179385.0,
    -1275806237668.0 / 842570457699.0])

    B = np.array([1432997174477.0 / 9575080441755.0,
    5161836677717.0 / 13612068292357.0,
    1720146321549.0 / 2090206949498.0,
    3134564353537.0 / 4481467310338.0,
    2277821191437.0 / 14882151754819.0])

    # --- parâmetros físicos ---
    eps = 8.854187817e-12  # Vacuum permittivity (F/m)
    mu = 4 * np.pi * 1e-7  # Vacuum permeability (H/m)

    # --- domínio ---
    vp = 1.0 * -1
    kx = 2 * np.pi
    xl, xr = 0, 1

    # --- retorna matriz D - diferenciação ---
    def get_D_matrix(nodes):
        N = len(nodes) - 1
        w = np.ones(N + 1)
        
        for j in range(N + 1):
            for i in range(N + 1):
                if i != j:
                    w[j] *= 1.0 / (nodes[j] - nodes[i])
        D = np.zeros((N + 1, N + 1))
        
        for i in range(N + 1):
            for j in range(N + 1):
                if i != j:
                    D[i, j] = (w[j] / w[i]) / (nodes[i] - nodes[j])
            D[i, i] = -np.sum(D[i, :])
            
        return D

    # --- retorna fluxo numérico ---
    def numerical_flux(uL, uR, a, flux_type):
        match flux_type:
            case 0:
                return 0.5 * a * (uL + uR)
            case 1:
                return a * uL if a > 0 else a * uR
            case _: # fallback central flux
                return 0.5 * a * (uL + uR)
            
    # --- retorna pontos GLL de ordem polinomial N ---
    def get_GLL_nodes(N):
        internal_nodes, _ = scipy.special.roots_jacobi(N - 1, 1, 1)
        r = np.concatenate([[-1.0], internal_nodes, [1.0]]) # elemento de referência
        return r

    # --- retorna pesos dos nós GLL de ordem polinomial N ---
    def get_GLL_weights(nodes, N):
        N = len(nodes) - 1

        # Legendre polynomial of degree N
        P = scipy.special.legendre(N)

        # Evaluate P_N(x_i)
        PN = np.polyval(P, nodes)

        # Compute weights
        w = 2.0 / (N * (N + 1) * PN**2)

        return w
        
    # --- operador DG ---
    def build_dg_operator(K, N, a, flux_type):
        dx_elem = (xr - xl) / K
        xi = get_GLL_nodes(N)
        D_ref = get_D_matrix(xi)
        D_phys = (2.0 / dx_elem) * D_ref
        w = get_GLL_weights(xi, N)
        M_local = w * dx_elem / 2.0
        
        # Coordenadas físicas dos nós
        x_nodes = np.zeros((K, N + 1))
        for k in range(K):
            xL = xl + k * dx_elem
            x_nodes[k, :] = xL + 0.5 * (xi + 1) * dx_elem
            
        return {
            'K': K, 'N': N, 'dx_elem': dx_elem,
            'D_phys': D_phys, 'M_local': M_local,
            'x_nodes': x_nodes, 'a': a, 'flux_type': flux_type
        }
        
    def rhs(u_flat, op):
        # puxa os dados do operador DG
        K, N = op['K'], op['N']
        D = op['D_phys']
        M = op['M_local']
        a = op['a']
        flux_type = op['flux_type']
        
        u = u_flat.reshape((K, N + 1))
        du = np.zeros_like(u)
        
        # itera sobre cada elemento e vai resolvendo em cada um
        for k in range(K):
            du[k, :] = -a * (D @ u[k, :])
            uL = u[k, N]
            uR = u[(k + 1) % K, 0]
            f_estrela_R = numerical_flux(uL, uR, a, flux_type)
            du[k, N] += a * (uL - f_estrela_R / a) / M[N] 
            
            uL_prev = u[(k - 1) % K, N]
            uR_curr = u[k, 0]
            f_estrela_L = numerical_flux(uL_prev, uR_curr, a, flux_type)
            du[k, 0] += a * (-uR_curr + f_estrela_L / a) / M[0]
            
        return du.flatten()

    # low storage Range Kutta quarta ordem
    def lsrk4_step(u, rhs_func, dt, op):
        q = np.zeros_like(u)
        for i in range(5):
            q = A[i] * q + dt * rhs_func(u, op)
            u = u + B[i] * q
            
        return u

    def DGTD(K, N, CFL, T_final, a, kx, flux_type):
        # gera o operador
        dg_operator = build_dg_operator(K, N, a, flux_type)
        x_nodes = dg_operator['x_nodes']
        
        # Condição inicial
        u = np.sin(kx * x_nodes).flatten().copy()
        u = np.exp(-0.5 * ((x_nodes - 0.5) / 0.05)**2).flatten().copy()
        
        # discretização
        dx = dg_operator['dx_elem']
        dt = CFL * dx / (np.abs(a) * (N + 1) ** 2)
        Nt = max(1, int(T_final / dt))
        dt = T_final / Nt # evita erro de arrendondamento
        
        # snapshots
        u_history = []
        snapshots_list = np.linspace(0, Nt - 10, 6)
        
        print("INICIANDO MARCHA NO TEMPO")
        print(f"dx = {dx:.3e}, dt = {dt:.3e}, Nt = {Nt}")
        
        # marcha no tempo
        for step in range(Nt):
            u = lsrk4_step(u, rhs, dt, dg_operator)
            
            if step in snapshots_list:
                u_history.append(u)
            
        return u_history, dx, dt, snapshots_list

    # --- Parâmetros de simulação ----
    K = 64 # numero de elementos
    N = 4 # ordem polinomial dentro dos elementos
    CFL = 0.2
    T_final = 1

    u_history, dx, dt, snapshots_list = DGTD(K, N, CFL, T_final, vp, kx, UPWIND_FLUX)
    \end{lstlisting}

    \chapter{Código Exercício 1.8}\label{cap:ex1-8-codigo}
    
    \begin{lstlisting}[language=Python, breaklines=true, basicstyle=\tiny]
    import numpy as np
    import matplotlib.pyplot as plt
    import scipy.special

    # --- constantes ----
    CENTRAL_FLUX = 0
    UPWIND_FLUX = 1

    A = np.array([0.0,
    -567301805773.0 / 1357537059087.0,
    -2404267990393.0 / 2016746695238.0,
    -3550918686646.0 / 2091501179385.0,
    -1275806237668.0 / 842570457699.0])

    B = np.array([1432997174477.0 / 9575080441755.0,
    5161836677717.0 / 13612068292357.0,
    1720146321549.0 / 2090206949498.0,
    3134564353537.0 / 4481467310338.0,
    2277821191437.0 / 14882151754819.0])

    # --- parâmetros físicos ---
    eps = 8.854187817e-12  # Vacuum permittivity (F/m)
    mu = 4 * np.pi * 1e-7  # Vacuum permeability (H/m)

    # --- domínio ---
    vp = 2 * np.pi
    kx = 2 * np.pi
    xl, xr = 0, 4 * np.pi

    # --- retorna matriz D - diferenciação ---
    def get_D_matrix(nodes):
        N = len(nodes) - 1
        w = np.ones(N + 1)
        
        for j in range(N + 1):
            for i in range(N + 1):
                if i != j:
                    w[j] *= 1.0 / (nodes[j] - nodes[i])
        D = np.zeros((N + 1, N + 1))
        
        for i in range(N + 1):
            for j in range(N + 1):
                if i != j:
                    D[i, j] = (w[j] / w[i]) / (nodes[i] - nodes[j])
            D[i, i] = -np.sum(D[i, :])
            
        return D

    # --- retorna fluxo numérico ---
    def numerical_flux(uL, uR, a, flux_type):
        match flux_type:
            case 0:
                return 0.5 * a * (uL + uR)
            case 1:
                return a * uL if a > 0 else a * uR
            case _: # fallback central flux
                return 0.5 * a * (uL + uR)
            
    # --- retorna pontos GLL de ordem polinomial N ---
    def get_GLL_nodes(N):
        internal_nodes, _ = scipy.special.roots_jacobi(N - 1, 1, 1)
        r = np.concatenate([[-1.0], internal_nodes, [1.0]]) # elemento de referência
        return r

    # --- retorna pesos dos nós GLL de ordem polinomial N ---
    def get_GLL_weights(nodes, N):
        N = len(nodes) - 1

        # Legendre polynomial of degree N
        P = scipy.special.legendre(N)

        # Evaluate P_N(x_i)
        PN = np.polyval(P, nodes)

        # Compute weights
        w = 2.0 / (N * (N + 1) * PN**2)

        return w
        
    # --- operador DG ---
    def build_dg_operator(K, N, a, flux_type):
        dx_elem = (xr - xl) / K
        xi = get_GLL_nodes(N)
        D_ref = get_D_matrix(xi)
        D_phys = (2.0 / dx_elem) * D_ref
        w = get_GLL_weights(xi, N)
        M_local = w * dx_elem / 2.0
        
        # Coordenadas físicas dos nós
        x_nodes = np.zeros((K, N + 1))
        for k in range(K):
            xL = xl + k * dx_elem
            x_nodes[k, :] = xL + 0.5 * (xi + 1) * dx_elem
            
        return {
            'K': K, 'N': N, 'dx_elem': dx_elem,
            'D_phys': D_phys, 'M_local': M_local,
            'x_nodes': x_nodes, 'a': a, 'flux_type': flux_type
        }
        
    def rhs(u_flat, op, step, dt):
        # puxa os dados do operador DG
        K, N = op['K'], op['N']
        D = op['D_phys']
        M = op['M_local']
        a = op['a']
        flux_type = op['flux_type']
        
        u = u_flat.reshape((K, N + 1))
        du = np.zeros_like(u)
        
        # itera sobre cada elemento e vai resolvendo em cada um
        for k in range(K):
            du[k, :] = -a * (D @ u[k, :])
            
            # fronteira direita
            if k == K - 1:
                uL = u[k, N]
                uR = uL
            else:
                uL = u[k, N]
                uR = u[k+1, 0]
            f_estrela_R = numerical_flux(uL, uR, a, flux_type)
            du[k, N] += a * (uL - f_estrela_R / a) / M[N]
            
            # fronteira esquerda
            if k == 0:
                uL = -np.sin(2.0 * np.pi * step * dt) # fonte
                uR = u[k, 0]
            else:
                uL = u[k-1, N]
                uR = u[k, 0]
            f_estrela_L = numerical_flux(uL, uR, a, flux_type)
            du[k, 0] += a * (-uR + f_estrela_L / a) / M[0]
            
        return du.flatten()

    # low storage Range Kutta quarta ordem
    def lsrk4_step(u, rhs_func, dt, op, step):
        q = np.zeros_like(u)
        for i in range(5):
            q = A[i] * q + dt * rhs_func(u, op, step, dt)
            u = u + B[i] * q
            
        return u

    def DGTD(K, N, CFL, T_final, a, kx, flux_type):
        # gera o operador
        dg_operator = build_dg_operator(K, N, a, flux_type)
        x_nodes = dg_operator['x_nodes']
        
        # Condição inicial
        u = np.sin(kx * x_nodes).flatten().copy()
        u = np.exp(-0.5 * ((x_nodes - 0.5) / 0.05)**2).flatten().copy()
        u = np.zeros_like(x_nodes.flatten())
        
        # discretização
        dx = dg_operator['dx_elem']
        dt = CFL * dx / (np.abs(a) * (N + 1) ** 2)
        Nt = max(1, int(T_final / dt))
        dt = T_final / Nt # evita erro de arrendondamento
        
        # snapshots
        u_history = []
        snapshots_list = np.linspace(0, Nt - 10, 6)
        
        print("INICIANDO MARCHA NO TEMPO")
        print(f"dx = {dx:.3e}, dt = {dt:.3e}, Nt = {Nt}")
        
        # marcha no tempo
        for step in range(Nt):
            u = lsrk4_step(u, rhs, dt, dg_operator, step)
            
            if step in snapshots_list:
                u_history.append(u)
            
        return u_history, dx, dt, snapshots_list

    # --- Parâmetros de simulação ----
    K = 32 # numero de elementos
    N = 4 # ordem polinomial dentro dos elementos
    CFL = 0.2
    T_final = 10

    u_history, dx, dt, snapshots_list = DGTD(K, N, CFL, T_final, vp, kx, UPWIND_FLUX)
    \end{lstlisting}

\end{anexosenv}
